\section{Planned Evaluation}
\label{sec:evaluation}

The evaluation should first establish whether the system behaves
correctly, then examine whether its guidance is useful. Three questions
provide a starting protocol.

\paragraph{RQ1: Is the instructional content correct and traceable?}
Subject reviewers should assess sampled explanations, question keys,
feedback, and source links against the original materials and explicit
rubrics. The audit should report the sampling procedure, number of
reviewed items, error categories, and reviewer agreement. Generated
examples should be checked separately from source quotations.

\paragraph{RQ2: Does the workflow support independent study and continuity?}
A formative study should observe students completing a topic, requesting
help, pausing, and returning to their work. Relevant outcomes include
task completion, recovery after interruption, and whether students
understand the distinction between assisted and independent performance.
A later controlled comparison could use the same course materials,
model, and study-time budget in StudyBirds and a source-grounded chat
baseline. Independent post-study questions and delayed assessment would
be needed to investigate learning and retention.

\paragraph{RQ3: Can instructors act on the reported difficulties?}
Instructors should review source-linked summaries and judge their
correctness, specificity, and usefulness for a teaching decision.
The analysis should also measure unsupported diagnoses and the effort
required to verify a summary.

Participant recruitment, study size, course selection, and the comparison
protocol remain to be specified. A reproducible evaluation should record
the system revision, model and prompt configuration, content versions,
latency, and operating cost. No participant outcomes or comparative
results are reported in this draft.
